\section{The all-place reversibility question}\label{sec:introduction}

A polynomial automorphism can be conjugate to its inverse over every
completion of a number field without being conjugate to its inverse
over that field. We give an explicit example in the group of polynomial
automorphisms of the affine plane. The distinction concerns existence
of a polynomial reversor, not just existence of a reversor in a selected
affine family.

For a field $L$, let $\Aut_L(\A^2)$ denote the group of polynomial maps
of the affine plane with polynomial inverses, all defined over $L$.
Group multiplication is composition, with the rightmost map acting
first. For $F\in\Aut_L(\A^2)$, put
\[
  \Rev_L(F)=
  \{R\in\Aut_L(\A^2):RFR^{-1}=F^{-1}\}.
\]
We call $F$ \emph{reversible over $L$} when this set is nonempty.
Our notation $\Rev_L(F)$ includes reversors only, not the union of
reversors and commuting symmetries.

Let $\Omega_K$ be the set of all finite and infinite places of a
number field $K$. The local-global question is whether
\begin{equation}\label{eq:hasse-question}
  \bigl[\Rev_{K_v}(F)\ne\varnothing
       \text{ for every }v\in\Omega_K\bigr]
  \quad\Longrightarrow\quad
  \Rev_K(F)\ne\varnothing
\end{equation}
for every finite nonempty composition of generalized H\'enon maps
\[
  H_i(x,y)=(y,P_i(y)-a_i x),\qquad
  P_i\in K[y],\quad \deg P_i\ge2,\quad a_i\in K^\times.
\]
There is no degree or order bound on a proposed reversor. No completion
may be enlarged or omitted, and $F$ is the full ordered composition,
not an iterate chosen in its place.

For $c\ne0$ and $d\ge2$, write
\begin{equation}\label{eq:hcd}
  H_{c,d}(x,y)=(y,cy^d-x).
\end{equation}
Its inverse is $(x,y)\mapsto(cx^d-y,x)$.
Each such map therefore satisfies the preceding H\'enon hypotheses.

\begin{theorem}\label{thm:counterexample}
Set
\begin{equation}\label{eq:counterexample}
  K=\Q(\sqrt7),\qquad
  F=H_{4,7}H_{1/16,15}H_{16,15}H_{1/4,7}.
\end{equation}
For every $v\in\Omega_K$ there is an affine involution
$R_v\in\Aut_{K_v}(\A^2)$ such that $R_vFR_v^{-1}=F^{-1}$.
Nevertheless,
\[
  \Rev_K(F)=\varnothing.
\]
In particular, implication \eqref{eq:hasse-question} is false even when
the search for a global reversor includes all polynomial degrees and
orders.
\end{theorem}

Although $F$ has rational coefficients, the base field in
Theorem~\ref{thm:counterexample} is the specified field $K$.
Both of its infinite places are real; it has no complex places.
The local witnesses and global exclusion will be proved over precisely
these fields.

The main algebraic step is an exhaustive calculation for a scalar
family. For a number field $L$ and $a\in L^\times$, define
\begin{equation}\label{eq:family}
  F_a=H_{a,7}H_{a^{-2},15}H_{a^2,15}H_{a^{-1},7},
  \qquad R_t(x,y)=(t^{-1}y,tx)\quad(t\ne0).
\end{equation}

\begin{theorem}\label{thm:family}
For every number field $L$ and every $a\in L^\times$,
\begin{equation}\label{eq:family-classification}
  \Rev_L(F_a)=
  \{F_a^jR_t:j\in\Z,\ t\in L^\times,\ t^8=a^2\}.
\end{equation}
Consequently, $F_a$ is reversible over $L$ if and only if
$a^2\in L^{\times8}$.
\end{theorem}

The integer $j$ in \eqref{eq:family-classification} is unrestricted.
The theorem is an equality for the entire polynomial reversor set,
not an affine ansatz. It will also provide the explicit local
construction: whenever $t^8=a^2$ in a characteristic-zero field,
$R_t$ reverses $F_a$.

\subsection*{Classical inputs and the distinction from conjugacy}

The polynomial group structure used here is classical.
The Jung--van der Kulk amalgam, its reduced words, and the associated
symmetry and reversor methods are described in
G\'omez--Meiss \cite{gomez2004reversors} and
Baake--Roberts \cite{baake2005symmetries}.
We specialize that machinery to one four-factor word and prove the
required full centralizer equality directly. A commuting-subgroup
inclusion would not suffice. In particular, the field hypothesis on
roots of unity in \cite[Theorem~2 and Corollary~1]{baake2005symmetries}
cannot be imposed during a calculation over $\C$.

The arithmetic obstruction is also classical. The exceptional
eighth-power class over $\Q(\sqrt7)$ is recorded in Song Wang's
account of the Grunwald--Wang theorem
\cite{wang2015grunwald}; we verify every place directly below.
Scalar twisting together with exhaustive centralizer descent is not
a new method: the final example of
Cantat--Dujardin \cite[Section~2.2]{cantat2024holomorphically}
already controls the field of polynomial conjugacy of two
scalar-twisted H\'enon maps through their full centralizer.
Combining that example with the Wang class gives an analogous
all-place failure for two independently specified maps.
That consequence of the cited mechanisms does not identify the
second map with the first map's inverse.
Here the remaining task is to realize the obstruction for the
constrained pair $(F,F^{-1})$ and exclude every alternative reversor.

The proof first determines the entire geometric centralizer, including
its translation image on the amalgam tree
(Sections~\ref{sec:axis}--\ref{sec:centralizer}).
Section~\ref{sec:descent} then proves
Theorem~\ref{thm:family} by a scalar twist and original-field descent.
Section~\ref{sec:places} proves Theorem~\ref{thm:counterexample}.
The two-factor control in Section~\ref{sec:control} explains why
failure of a selected affine swap cannot replace the full-group
argument.
